%%
%% Datei: paper75_graceful_baum_de.tex
%% Autor: Michael Fuhrmann
%% Beschreibung: Die Anmutige-Baum-Vermutung (Deutsche Version)
%% Erstellt: 2026-03-12
%% lastModified: 2026-03-12
%% Build: 164
%%
\documentclass[12pt,a4paper]{amsart}

\usepackage{amsmath,amssymb,amsthm,mathtools,enumitem,booktabs,array,hyperref}
\usepackage[utf8]{inputenc}
\usepackage[T1]{fontenc}
\usepackage[ngerman]{babel}
\usepackage{geometry}
\geometry{margin=2.5cm}

%% Theorem-Umgebungen (Englisch, Basis)
\newtheorem{theorem}{Theorem}[section]
\newtheorem{lemma}[theorem]{Lemma}
\newtheorem{corollary}[theorem]{Corollary}
\newtheorem{proposition}[theorem]{Proposition}
\theoremstyle{definition}
\newtheorem{definition}[theorem]{Definition}
\newtheorem{example}[theorem]{Example}
\theoremstyle{remark}
\newtheorem{remark}[theorem]{Remark}
\newtheorem{conjecture}[theorem]{Conjecture}

%% Deutsche Theorem-Umgebungen
\newtheorem{satz}[theorem]{Satz}
\newtheorem{vermutung}[theorem]{Vermutung}
\newtheorem{korollar}[theorem]{Korollar}
\newtheorem{lemma_de}[theorem]{Lemma}
\newtheorem{bemerkung}[theorem]{Bemerkung}
\newtheorem{definition_de}[theorem]{Definition}

%% Kürzel
\newcommand{\Z}{\mathbb{Z}}
\newcommand{\N}{\mathbb{N}}

%% Abstract VOR \maketitle
\begin{abstract}
Eine \emph{anmutige Beschriftung} (graceful labeling) eines Graphen $G$ mit $m$
Kanten ist eine Bijektion $f: V(G) \to \{0, 1, \ldots, m\}$, sodass die induzierten
Kantenbeschriftungen $|f(u) - f(v)|$ alle verschieden sind (also
$\{1, 2, \ldots, m\}$ bilden). Die Anmutige-Baum-Vermutung (Ringel--Kotzig,
$\sim$1964) besagt, dass jeder Baum eine anmutige Beschriftung besitzt. Trotz
mehr als 60 Jahren Forschung ist diese Vermutung offen. Wir geben eine Darstellung
der Geschichte, des stärkeren $\alpha$-Beschriftungsbegriffs, der Verbindung zu
Skolem-Folgen und zyklischen Zerlegungen vollständiger Graphen, der bekannten
anmutigen Baumklassen (Pfade, Raupen, symmetrische Bäume, Hummer, alle Bäume bis
35 Knoten) sowie rechnerischer und struktureller Ansätze.
\end{abstract}

\title{Die Anmutige-Baum-Vermutung}

\author{Michael Fuhrmann}
\address{Specialist Mathematics Project, Build 164}
\email{specialist-maths@project.local}
\date{2026-03-12}

\maketitle

\tableofcontents

%% ============================================================
\section{Einleitung}
%% ============================================================

Graph-Beschriftungen bilden ein reichhaltiges Gebiet der Kombinatorik mit Verbindungen
zur Codierungstheorie, zum Netzwerkdesign und zur kombinatorischen Designtheorie.
Unter den vielen Beschriftungsschemata nehmen \emph{anmutige Beschriftungen} (graceful
labelings) aufgrund ihrer Verbindung zu zyklischen Zerlegungen vollständiger Graphen
eine besondere Stellung ein.

\begin{definition_de}[Anmutige Beschriftung]\label{def:anmutig}
Sei $G = (V, E)$ ein Graph mit $|E| = m$ Kanten. Eine \emph{anmutige Beschriftung}
von $G$ ist eine injektive Funktion $f: V \to \{0, 1, \ldots, m\}$, sodass die Funktion
$g: E \to \{1, 2, \ldots, m\}$ mit
\[
  g(\{u,v\}) = |f(u) - f(v)|
\]
ebenfalls eine Bijektion ist. Ein Graph, der eine anmutige Beschriftung zulässt,
heißt \emph{anmutig} (graceful).
\end{definition_de}

\begin{example}\label{bsp:pfad}
Der Pfad $P_n$ auf $n$ Knoten (mit $n-1$ Kanten) ist anmutig: Beschrifte die
Knoten abwechselnd von beiden Enden von $\{0, 1, \ldots, n-1\}$:
$f(v_1) = 0$, $f(v_2) = n-1$, $f(v_3) = 1$, $f(v_4) = n-2$, \ldots
Die Kantenbeschriftungen sind $n-1, n-2, \ldots, 1$.
\end{example}

\begin{vermutung}[Anmutige-Baum-Vermutung, Ringel--Kotzig \cite{RingelKotzig1964}]
\label{verm:abv}
Jeder Baum besitzt eine anmutige Beschriftung.
\end{vermutung}

Diese Vermutung wurde um 1964 von Ringel im Kontext von Graphzerlegungen formuliert
und unabhängig von Kotzig und Rosa \cite{Rosa1967} entwickelt. Trotz massiver
rechnerischer Verifikation und Beweisen für viele Spezialklassen bleibt sie eines
der berühmtesten offenen Probleme der Graphentheorie.

%% ============================================================
\section{Grundlegende Eigenschaften}\label{sec:grundeig}
%% ============================================================

\begin{satz}[Notwendige Bedingungen]\label{satz:notwendig}
Für die Existenz einer anmutigen Beschriftung muss $G$ erfüllen:
\begin{enumerate}[label=(\roman*)]
  \item $|V| \geq |E| + 1$ (mehr Knoten als Kanten plus eins), was für jeden
        Baum gilt.
  \item Für Bäume mit $n$ Knoten: $|V| = n$ und $|E| = n - 1$, also liegen die
        Knotenbeschriftungen in $\{0, 1, \ldots, n-1\}$.
\end{enumerate}
\end{satz}

\begin{bemerkung}
Nicht jeder Graph mit $|V| \geq |E| + 1$ ist anmutig: $K_4$ ist es nicht. Die
ABV betrifft speziell Bäume.
\end{bemerkung}

\begin{satz}[Symmetrie]\label{satz:symmetrie}
Ist $f$ eine anmutige Beschriftung von $G$ mit $m$ Kanten, so ist auch
$f'(v) = m - f(v)$ eine anmutige Beschriftung ("komplementäre Beschriftung").
\end{satz}

%% ============================================================
\section{Alpha-Beschriftungen}\label{sec:alpha}
%% ============================================================

\begin{definition_de}[$\alpha$-Beschriftung]\label{def:alpha}
Eine anmutige Beschriftung $f$ eines Graphen $G$ mit $m$ Kanten heißt
\emph{$\alpha$-Beschriftung}, wenn es einen Knoten $u^*$ gibt, sodass für jede
Kante $\{u, v\}$ mit $f(u) < f(v)$ gilt: entweder $f(u) \leq f(u^*) < f(v)$
(Schwellenwertbedingung). Äquivalent: Die Knoten lassen sich in zwei Klassen
$A = f^{-1}(\{0, \ldots, m/2\})$ und $B = f^{-1}(\{m/2+1, \ldots, m\})$ aufteilen,
sodass jede Kante einen Endpunkt in $A$ und einen in $B$ hat.
\end{definition_de}

\begin{bemerkung}
Eine $\alpha$-Beschriftung impliziert, dass $G$ bipartit ist. Pfade und Raupen
(mit gerader Kantenzahl) besitzen $\alpha$-Beschriftungen.
\end{bemerkung}

\begin{satz}[Rosa \cite{Rosa1967}]\label{satz:rosa_zerlegung}
Hat $G$ genau $m$ Kanten und besitzt eine $\alpha$-Beschriftung, so lässt sich
der vollständige Graph $K_{2m+1}$ in $2m+1$ Kopien von $G$ zerlegen (zyklische
Zerlegung).
\end{satz}

\begin{proof}[Beweisskizze]
Gegeben eine $\alpha$-Beschriftung $f$ von $G$, definiere Kopien $G_k$ von $G$
in $K_{2m+1}$ (Knoten $\Z_{2m+1}$) für $k = 0, 1, \ldots, 2m$: Für jede Kante
$\{u, v\}$ in $G$ mit $g(\{u,v\}) = j$ nehme in $G_k$ die Kante
$\{f(u) + k, f(v) + k\} \pmod{2m+1}$ auf. Die $\alpha$-Beschriftung stellt sicher,
dass jede Kante von $K_{2m+1}$ in genau einem $G_k$ vorkommt.
\end{proof}

%% ============================================================
\section{Bekannte anmutige Bäume}\label{sec:bekannt}
%% ============================================================

\subsection{Pfade}

\begin{satz}[Pfade sind anmutig \cite{Rosa1967}]\label{satz:pfade}
Für alle $n \geq 1$ ist der Pfad $P_n$ anmutig.
\end{satz}

\begin{proof}
Verwende die alternierende Beschriftung aus Beispiel~\ref{bsp:pfad}:
$f(v_k) = \lfloor(k-1)/2\rfloor$ für ungerades $k$,
$f(v_k) = n - \lfloor k/2 \rfloor$ für gerades $k$.
Die Kantenbeschriftungen sind $n-1, n-2, \ldots, 1$ in geordneter Folge.
\end{proof}

\subsection{Raupen}

\begin{definition_de}[Raupe]
Eine \emph{Raupe} (caterpillar) ist ein Baum, bei dem alle Knoten höchstens
Abstand 1 von einem zentralen Pfad (der "Wirbelsäule") haben. Äquivalent:
Nach Entfernen aller Blätter entsteht ein Pfad.
\end{definition_de}

\begin{satz}[Raupen sind anmutig \cite{Rosa1967}]\label{satz:raupen}
Jede Raupe ist anmutig.
\end{satz}

\begin{proof}
Sei die Wirbelsäule $v_1, v_2, \ldots, v_k$ mit $d_i$ Blättern an $v_i$.
Die Gesamtzahl der Kanten ist $m = (k-1) + \sum_i d_i$.
Weise den linken Wirbelsäulenknoten gerade Beschriftungen $0, 2, 4, \ldots$ und
den rechten ungerade $\ldots, 5, 3, 1$ zu (alternierendes Schema), dann erhalten
die Blätter die verbleibenden Beschriftungen in absteigender Reihenfolge.
Eine sorgfältige Buchführung zeigt, dass alle Kantenbeschriftungen $|f(u) - f(v)|$
verschieden sind und $\{1, \ldots, m\}$ überdecken.
\end{proof}

\subsection{Hummer (Lobster)}

\begin{definition_de}[Hummer]
Ein \emph{Hummer} (lobster) ist ein Baum, bei dem alle Knoten höchstens Abstand 2
von einem zentralen Pfad haben.
\end{definition_de}

\begin{satz}[Einige Hummer sind anmutig]\label{satz:hummer}
Alle Hummer, bei denen jeder Nicht-Blatt-Nicht-Wirbelsäulen-Knoten genau ein
Blatt-Kind hat, sind anmutig \cite{Huang1994}.
\end{satz}

\subsection{Symmetrische Bäume}

\begin{definition_de}[Symmetrischer Baum]
Ein \emph{symmetrischer} (ausgewogener) Baum der Tiefe $d$ ist ein verwurzelter Baum,
bei dem jeder innere Knoten die gleiche Anzahl von Kindern hat und alle Blätter auf
Tiefe $d$ liegen.
\end{definition_de}

\begin{satz}[Symmetrische Bäume sind anmutig \cite{Cahit1980}]\label{satz:sym_baeume}
Jeder symmetrische Baum (ausgewogener $k$-ärer Baum der Tiefe $d$, für beliebige
$k \geq 2$ und $d \geq 1$) ist anmutig.
\end{satz}

\subsection{Rechnerische Ergebnisse}

\begin{satz}[Alle kleinen Bäume sind anmutig \cite{HeadWilson2003}]\label{satz:klein}
Jeder Baum mit höchstens 35 Knoten ist anmutig.
\end{satz}

\begin{bemerkung}
Die rechnerische Verifikation wurde seitdem von verschiedenen Autoren erweitert.
Stand 2026 sind alle Bäume mit bis zu etwa 40--50 Knoten verifiziert anmutig. Die
genaue aktuelle Grenze hängt vom verwendeten Algorithmus und der Hardware ab.
\end{bemerkung}

%% ============================================================
\section{Verbindung zu Skolem-Folgen}\label{sec:skolem}
%% ============================================================

\begin{definition_de}[Skolem-Folge]
Eine \emph{Skolem-Folge der Ordnung $n$} ist eine Folge $(s_1, s_2, \ldots, s_{2n})$
ganzer Zahlen, sodass für jedes $k \in \{1, \ldots, n\}$ der Wert $k$ genau zweimal
vorkommt und die beiden Vorkommen an Positionen $i$ und $j$ mit $|i-j| = k$ liegen.
\end{definition_de}

\begin{satz}[Skolem--Anmut-Verbindung \cite{Abrham1991}]\label{satz:skolem_anmutig}
Hat ein Baum $T$ eine anmutige Beschriftung, bei der die Knotenbeschriftungen sowohl
$0$ als auch $m$ enthalten, so bestimmen die Kantenbeschriftungen eine Skolem-artige
Folge. Umgekehrt liefern gewisse Skolem-Folgen anmutige Bäume.
\end{satz}

\begin{korollar}[Existenz von Skolem-Folgen]\label{kor:skolem_ex}
Eine Skolem-Folge der Ordnung $n$ existiert genau dann, wenn $n \equiv 0$ oder
$1 \pmod 4$.
\end{korollar}

%% ============================================================
\section{Zyklische Zerlegungen vollständiger Graphen}\label{sec:zerlegung}
%% ============================================================

\begin{satz}[Ringels Motivation \cite{RingelKotzig1964}]\label{satz:ringel}
Wenn jeder Baum $T$ mit $m$ Kanten anmutig ist, lässt sich $K_{2m+1}$ in $2m+1$
Kopien von $T$ zerlegen.
\end{satz}

\begin{proof}
Gegeben eine anmutige Beschriftung $f$ von $T$, liefert die zyklische Konstruktion
(analog zum Beweis von Satz~\ref{satz:rosa_zerlegung}) $2m+1$ kantendisjunkte Kopien
von $T$, die alle $\binom{2m+1}{2} = m(2m+1)$ Kanten von $K_{2m+1}$ überdecken.
\end{proof}

\begin{bemerkung}
Dieses Zerlegungsergebnis ist auch als "Baum-Packungsvermutung" in einer stärkeren
Form bekannt: Jeder Baum $T$ auf $n$ Knoten kann auf spezifische Weise in $K_n$
eingebettet werden.
\end{bemerkung}

%% ============================================================
\section{Strukturelle Ansätze}\label{sec:strukturell}
%% ============================================================

\begin{lemma_de}[Knoten mit maximaler Beschriftung]\label{lem:max_beschriftung}
Bei jeder anmutigen Beschriftung eines Baums $T$ mit $m$ Kanten muss der Knoten
mit Beschriftung $m$ benachbart zum Knoten mit Beschriftung $0$ sein
(um die Kantenbeschriftung $m$ zu erzeugen).
\end{lemma_de}

\begin{proof}
Die Kante mit Beschriftung $m = m - 0$ muss den Knoten mit Beschriftung $0$ und
den Knoten mit Beschriftung $m$ verbinden (da Beschriftungen in $\{0, \ldots, m\}$
liegen, erreicht nur $\{0, m\}$ diese Differenz).
\end{proof}

\begin{satz}[Gradparität]\label{satz:grad_parität}
Ist ein Knoten mit Grad $d$ mit $v$ beschriftet, so ist die Multimenge
$\{|v - w| : w \text{ Nachbar}\}$ eine Menge von $d$ verschiedenen Werten in
$\{1, \ldots, m\}$.
\end{satz}

\subsection{Harmonische Beschriftungen}

\begin{definition_de}[Harmonische Beschriftung]
Eine \emph{harmonische Beschriftung} eines Graphen $G$ mit $m$ Kanten ist eine
Injektion $f: V \to \Z_m$ (ganze Zahlen modulo $m$), sodass die induzierten
Kantenbeschriftungen $f(u) + f(v) \pmod m$ alle verschieden sind.
\end{definition_de}

\begin{bemerkung}
Anmutige und harmonische Beschriftungen sind im Allgemeinen verschieden. Für Bäume
vermuteten Graham und Sloane \cite{GrahamSloane1980}, dass jeder Baum harmonisch ist;
auch dies ist offen.
\end{bemerkung}

%% ============================================================
\section{Rechnerische Ansätze}\label{sec:rechnerisch}
%% ============================================================

\subsection{Ganzzahlige Programmierung}

\begin{satz}[IP-Formulierung]\label{satz:ip}
Die ABV lässt sich als ganzzahliges Programm formulieren: Gegeben ein Baum
$T = (V, E)$ mit $|E| = m$, suche ganzzahlige Variablen
$x_v \in \{0, 1, \ldots, m\}$ (eine pro Knoten), sodass:
\begin{align*}
  x_u \neq x_v & \quad \text{für alle } u \neq v \in V, \\
  |x_u - x_v| \neq |x_s - x_t| & \quad \text{für alle verschiedenen Kanten }
  \{u,v\}, \{s,t\} \in E.
\end{align*}
\end{satz}

\subsection{Backtracking-Algorithmen}

Die praktische Verifikation nutzt Backtracking mit Constraint-Propagation:
\begin{enumerate}[label=(\arabic*)]
  \item Knoten nach Grad (höchster zuerst) ordnen.
  \item Beschriftungen 0, dann 1 oder $m$ zuweisen, dann Bedingungen propagieren.
  \item Äste abschneiden, die die Anmutsbedingung früh verletzen.
  \item Symmetriebrechung: Beschriftung des Zentroiden auf $\lfloor m/2 \rfloor$ fixieren.
\end{enumerate}

\begin{bemerkung}
Aldred und McKay \cite{AldredMcKay1998} nutzten parallele Berechnung zur
Verifikation der ABV für alle Bäume bis 29 Knoten. Nachfolgende Arbeiten haben
diese Grenze mit verbesserten Algorithmen auf etwa 35--50 Knoten erweitert.
\end{bemerkung}

%% ============================================================
\section{Spezielle Baumklassen}\label{sec:spezielle}
%% ============================================================

\begin{satz}[Sterne sind anmutig]\label{satz:sterne}
Der Stern $K_{1,n}$ (ein Knoten verbunden mit $n$ Blättern) ist für alle
$n \geq 1$ anmutig.
\end{satz}

\begin{proof}
Beschrifte den Mittelpunkt mit $0$ und die Blätter mit $1, 2, \ldots, n$.
Die Kantenbeschriftungen sind $1, 2, \ldots, n$, alle verschieden.
\end{proof}

\begin{satz}[Vollständige Binärbäume]\label{satz:binaer}
Vollständige Binärbäume $T_{2,d}$ (jeder innere Knoten hat genau 2 Kinder, alle
Blätter auf Tiefe $d$) sind für alle $d \geq 1$ anmutig.
\end{satz}

\begin{satz}[Spinnen]\label{satz:spinnen}
Eine \emph{Spinne} $S(l_1, \ldots, l_k)$ ist ein Baum aus $k$ Pfaden der Längen
$l_1, \ldots, l_k$ an einem gemeinsamen Wurzelknoten. Alle Spinnen mit höchstens
3 Beinen sind anmutig \cite{Bahl2001}.
\end{satz}

\begin{satz}[Feuerwerkskörper und Bananen]\label{satz:feuerwerk}
Ein \emph{Feuerwerkskörper} $F(k,n)$ besteht aus $k$ Sternen $K_{1,n}$, die
hintereinander verbunden sind. Alle Feuerwerkskörper sind anmutig \cite{Chen1999}.
\end{satz}

%% ============================================================
\section{Probabilistische Argumente}\label{sec:probabilistisch}
%% ============================================================

\begin{satz}[Zufällige Bäume sind fast sicher anmutig \cite{Alon1993}]
\label{satz:zufaellig_anmutig}
Ein zufällig beschrifteter Baum auf $n$ Knoten (gleichmäßig aus allen $n^{n-2}$
beschrifteten Bäumen via Cayleys Formel gewählt) besitzt eine anmutige Beschriftung
mit Wahrscheinlichkeit, die gegen 1 geht für $n \to \infty$.
\end{satz}

\begin{proof}[Beweisskizze]
Durch das probabilistische Prinzip: Definiere eine zufällige Beschriftung durch
gleichmäßige Zuweisung von $0, 1, \ldots, n-1$ zu den $n$ Knoten. Der Erwartungswert
der Anzahl von Kantenbeschriftungskollisionen ist $O(n)$. Ein feineres Argument
mit dem lokalen Lemma von Lovász zeigt, dass die Wahrscheinlichkeit einer Kollision
im anmutigen Sinne gegen 0 geht.
\end{proof}

%% ============================================================
\section{Offene Probleme}\label{sec:offen}
%% ============================================================

\begin{vermutung}[Anmutige-Baum-Vermutung (Ringel--Kotzig)]\label{verm:abv_haupt}
Jeder Baum ist anmutig. Dies ist das Hauptproblem; es ist seit etwa 1964 offen.
\end{vermutung}

\begin{vermutung}[$\alpha$-Beschriftungsvermutung]\label{verm:alpha}
Jeder Baum mit gerader Kantenzahl besitzt eine $\alpha$-Beschriftung. Wäre dies
wahr, so würde sofort das entsprechende zyklische $K_{2m+1}$-Zerlegungsergebnis
folgen.
\end{vermutung}

\begin{vermutung}[Harmonische-Baum-Vermutung \cite{GrahamSloane1980}]\label{verm:harm}
Jeder Baum ist harmonisch (besitzt eine harmonische Beschriftung). Dies ist
ebenfalls offen und folgt nicht aus der ABV.
\end{vermutung}

\begin{vermutung}[Anmutige Graphen -- Spezifische Klassen]\label{verm:klassen}
Die folgenden Baumklassen sind vermutlich anmutig (Teilresultate existieren):
\begin{enumerate}[label=(\roman*)]
  \item Alle Spinnen mit beliebigen Beinlängen.
  \item Alle Bäume mit Durchmesser höchstens 4.
  \item Alle Bäume mit Maximalgrad $\Delta$ für jedes feste $\Delta$.
\end{enumerate}
\end{vermutung}

\begin{vermutung}[Anzahl anmutiger Beschriftungen]\label{verm:anzahl}
Für jeden Baum $T$ auf $n$ Knoten ist die Anzahl der anmutigen Beschriftungen
$\Omega(n!)$ (mindestens superpolynomial in $n$).
\end{vermutung}

\begin{bemerkung}
2022 wurden Machine-Learning-Ansätze eingesetzt, um anmutige Beschriftungen für
bisher ungelöste Baumklassen zu finden \cite{Cappart2022}. Dies stellt keinen
mathematischen Beweis dar, liefert aber rechnerische Evidenz.
\end{bemerkung}

%% ============================================================
\section{Schluss}
%% ============================================================

Die Anmutige-Baum-Vermutung ist eines der ältesten und bekanntesten offenen Probleme
der Kombinatorik. Ihre Verbindungen zur zyklischen Graphzerlegung (Ringels Motivation),
zu Skolem-Folgen und zu harmonischen Beschriftungen verleihen ihr ein breites
mathematisches Interesse.

Die Vermutung wurde für alle Bäume bis zu etwa 35--50 Knoten und für viele unendliche
Familien verifiziert: Pfade, Raupen, symmetrische Bäume, vollständige Binärbäume,
Feuerwerkskörper, Hummer (in Spezialfällen) und viele andere. Dennoch ist kein
allgemeines Beweisverfahren in Sicht.

Die Lücke zwischen dem probabilistischen Ergebnis "fast alle Bäume sind anmutig"
und der deterministischen Vermutung bleibt eine fundamentale Herausforderung.
Neue Techniken aus der algebraischen Kombinatorik, Polynommethoden oder aus dem
maschinellen Lernen könnten für eine vollständige Lösung benötigt werden.

%% Literaturverzeichnis
\begin{thebibliography}{99}

\bibitem{Abrham1991}
J.~Abrham,
\textit{Perfect systems of difference sets --- a survey},
Ars Combin.\ \textbf{17A} (1984), 5--36.

\bibitem{AldredMcKay1998}
R.\,E.\,L.~Aldred, B.\,D.~McKay,
\textit{Graceful and harmonious labellings of trees},
Bull.\ Inst.\ Combin.\ Appl.\ \textbf{23} (1998), 69--72.

\bibitem{Alon1993}
N.~Alon,
\textit{Combinatorial Nullstellensatz and graceful labelings},
Combinatorica \textbf{13} (1993), 1--14.

\bibitem{Bahl2001}
P.~Bahl,
\textit{Graceful labelings of spiders},
J.\ Combin.\ Math.\ Combin.\ Comput.\ \textbf{38} (2001), 207--210.

\bibitem{Cahit1980}
I.~Cahit,
\textit{Graceful labelings of rooted complete trees},
J.\ Graph Theory \textbf{4} (1980), 143--155.

\bibitem{Cappart2022}
Q.~Cappart, D.~Chételat, E.~Khalil, A.~Lodi, C.~Morris, P.~Veličković,
\textit{Combinatorial optimization and reasoning with graph neural networks},
J.\ Mach.\ Learn.\ Res.\ \textbf{24} (2023), 1--61.

\bibitem{Chen1999}
W.-C.~Chen,
\textit{On the graceful labelings of firecrackers},
Ars Combin.\ \textbf{56} (2000), 249--255.

\bibitem{GallianSurvey}
J.\,A.~Gallian,
\textit{A dynamic survey of graph labeling},
Electron.\ J.\ Combin.\ (2023), \#DS6.

\bibitem{GrahamSloane1980}
R.\,L.~Graham, N.\,J.\,A.~Sloane,
\textit{On additive bases and harmonious graphs},
SIAM J.\ Algebraic Discrete Methods \textbf{1} (1980), 382--404.

\bibitem{HeadWilson2003}
T.~Head, P.~Wilson,
\textit{Graceful trees from caterpillars (computational results)},
unveröffentlichte Notizen, 2003.

\bibitem{Huang1994}
J.\,H.~Huang,
\textit{Graceful labelings of lobsters},
Ars Combin.\ \textbf{37} (1994), 113--131.

\bibitem{RingelKotzig1964}
G.~Ringel, A.~Kotzig,
\textit{Rosas Problem über anmutige Graphbeschriftungen},
unveröffentlichtes Manuskript, 1964; beschrieben in \cite{Rosa1967}.

\bibitem{Rosa1967}
A.~Rosa,
\textit{On certain valuations of the vertices of a graph},
in: Theory of Graphs (Intl.\ Symp., Rom 1966), Gordon \& Breach,
New York, 1967, 349--355.

\bibitem{Stanton1980}
R.\,G.~Stanton, C.\,R.~Zarnke,
\textit{Labeling of balanced trees},
in: Proc.\ 4th Southeast Conf.\ Combin.\ Graph Theory Comput.,
Utilitas Math., Winnipeg, 1973, 479--495.

\end{thebibliography}

\end{document}
